% Lösungen Einheit 2 von 3 – Kathete und Umkehrung (zusammenbau v0.1)
\einheitenkopf{Einheit 2 von 3 · Kathete und Umkehrung}

\erg{17}{a) $8{,}5$ cm}
%% TODO Lösung der Erklärzeile 18g) fehlt in der Bank
\erg{18}{a) Kathete gesucht – minus \quad b) $a^2 = c^2 - b^2 = 7\,921 - 6\,400 = 1\,521$, $a = 39$ cm \quad c) $b^2 = c^2 - a^2 = 9\,409 - 4\,225 = 5\,184$, $b = 72$ m \quad d) $a^2 = c^2 - b^2 = 5\,329 - 2\,304 = 3\,025$, $a = 55$ mm \quad e) $b^2 = c^2 - a^2 = 7\,569 - 3\,600 = 3\,969$, $b = 63$ cm \quad f) $a^2 = c^2 - b^2 = 2\,500 - 196 = 2\,304$, $a = 48$ m \quad g) \ldots \quad h) $a^2 = 85$, $a \approx 9{,}2$ cm \quad i) $a^2 = 36$, $a = 6$ m \quad j) $a^2 = 36$, $a = 6$ cm; $6 < 7{,}5$, passt \quad k) $f = \sqrt{g^2 - e^2}$ (Hypotenusenquadrat minus Quadrat der anderen Kathete) \quad l) $\sqrt{8^2 - 4{,}9^2} = \sqrt{39{,}99} \approx 6{,}3$ cm, die Angabe stimmt \quad m) $250^2 - 240^2 = 4\,900$, Höhenunterschied $= 70$ m \quad n) $33^2 + 56^2 = 4\,225$ und $65^2 = 4\,225$, gleich – rechtwinklig \quad o) $1{,}8^2 + 2{,}4^2 = 9 = 3^2$; nach der Umkehrung des Satzes ist das Dreieck rechtwinklig, der rechte Winkel liegt gegenüber der $3{,}0$-m-Seite. \quad p) $h = \sqrt{426^2 - 390^2} = \sqrt{29\,376} \approx 171{,}4$ m \quad q) Teildreieck mit rechtem Winkel bei der Höhe; $PF = \sqrt{9{,}6^2 - 5{,}3^2} = \sqrt{64{,}07} \approx 8{,}0$ m \quad r) $AD^2 = 32^2 + 32^2 = 2\,048$, ebenso $DC^2 = 2\,048$; $AD^2 + DC^2 = 4\,096 = 64^2 = AC^2$, also nach der Umkehrung ein rechter Winkel bei D.}
\erg{19}{a) $11^2 + 60^2 = 3\,721 = 61^2$ – ja}
\erg{20}{a) Quadrate addiert, obwohl eine Kathete gesucht ist – das Ergebnis ist sogar länger als die Hypotenuse. Richtig: $a^2 = 58^2 - 40^2 = 1\,764$, $a = 42$ cm \quad b) Das Hypotenusenquadrat besteht aus beiden Kathetenquadraten zusammen; nimmt man das bekannte Kathetenquadrat weg, bleibt das gesuchte übrig. \quad c) $\sqrt{30^2 - 14^2} = \sqrt{704} \approx 26{,}5$ m; dazu die $3$ m: $26{,}5 + 3 \approx 29{,}5$ m – ja, sie reicht. \quad d) Hypotenuse $7{,}3$ gegenüber dem rechten Winkel, Katheten $4{,}8$ und $b$.}

20d)

\dreieckrw{4.36}{5}{$4{,}8$}{$b$}{$7{,}3$}

